\documentclass[12pt]{article}
\usepackage[utf8]{inputenc}
\usepackage[T1]{fontenc}
\usepackage{hyperref}
\title{Ten'edn Language Resources for K-12 Curriculum Development}
\author{Compiled by Claude Sonnet 5 (claude-sonnet-5) for Andr\'as Kornai}
\date{}
\begin{document}
\maketitle

\section{Introduction}
Ten'edn (ISO 639-3 tnz, also called Mos, Tonga or Tonga-Mos) is an Aslian (Austroasiatic) language of the Maniq people of southern Thailand and northern Malaysia. About 365 native speakers are reported (2007--2014). It is distinct from Kensiu (kns), which is also spoken by Maniq groups. Sources report that it has no standard writing system. No school materials were found.

\section{Primary Educational Resources}
None found.

\section{Linguistic and Reference Materials}
\begin{itemize}
\item Wikipedia article on Ten'edn: \url{https://en.wikipedia.org/wiki/Ten%27edn}
\item Max Planck EVA page on Ten'edn numerals: \url{https://lingweb.eva.mpg.de/channumerals/Ten'edn.htm}
\item Wikipedia article on the Maniq people: \url{https://wikipedia.com/wiki/Maniq}
\item Wikipedia article on Jahaic languages, for the wider Aslian context: \url{https://en.wikipedia.org/wiki/Jahaic_languages}
\end{itemize}

\section{Religious and Community Resources}
None found.

\section{Digital and Audio Resources}
None found.

\section{Curriculum Development Context}
The language is spoken by a small, mobile community and has no writing system, so any curriculum would first need community-led orthography and documentation work. Thai and Malay are the national school languages. The search returned only encyclopaedic sources.
\end{document}